\chapter{Implementación}
\label{cap:implementacion}

Este capítulo describe las herramientas, la organización del código y los
detalles de implementación del \textit{pipeline}.

\section{Herramientas y librerías}
El proyecto se desarrolla en \textbf{Python}. Las principales librerías son:

\begin{table}[H]
\centering
\caption{Principales herramientas utilizadas.}
\label{tab:herramientas}
\begin{tabular}{@{}ll@{}}
\toprule
\textbf{Ámbito} & \textbf{Herramienta} \\
\midrule
Manipulación de datos      & pandas, NumPy \\
Imagen                     & OpenCV, scikit-image, Pillow \\
TDA / homología persistente& giotto-tda, Ripser, GUDHI \\
Aprendizaje automático     & scikit-learn \\
Aprendizaje profundo       & TensorFlow / Keras (o PyTorch) \\
Visualización              & Matplotlib, seaborn \\
Entorno                    & Jupyter Notebook \\
\bottomrule
\end{tabular}
\end{table}

\section{Organización del código}
El repositorio se estructura de la siguiente forma (véase el
\texttt{README.md}):
\begin{lstlisting}[language=bash, numbers=none]
mammoviz-tda/
  data/            # datos (no versionados)
  src/             # codigo fuente del pipeline
    data.py        # carga y particion por paciente
    preprocessing.py  # preprocesado de imagenes
    topology.py    # homologia persistente y vectorizacion
    models.py      # clasificadores y CNN
    evaluate.py    # metricas y visualizacion
  notebooks/       # experimentos reproducibles
  tfg/             # memoria en LaTeX
\end{lstlisting}

\section{Extracción de descriptores topológicos}
El siguiente fragmento ilustra el cálculo de un diagrama de persistencia a partir
de una imagen en escala de grises mediante una filtración por \textit{sublevel
sets}, utilizando \texttt{giotto-tda}.

\begin{lstlisting}[language=Python, caption={Ejemplo de extracción de un diagrama de persistencia.}, label={lst:tda}]
import numpy as np
from gtda.images import HeightFiltration, RadialFiltration
from gtda.homology import CubicalPersistence
from gtda.diagrams import PersistenceImage

def descriptor_topologico(imagen_gris):
    # imagen_gris: array 2D normalizado en [0, 1]
    x = imagen_gris[None, :, :]              # (1, H, W)
    cubical = CubicalPersistence(homology_dimensions=(0, 1))
    diagramas = cubical.fit_transform(x)     # diagramas de persistencia
    pi = PersistenceImage(n_bins=20)
    vector = pi.fit_transform(diagramas)     # imagen de persistencia
    return vector.reshape(len(x), -1)
\end{lstlisting}

\section{Reproducibilidad}
Para garantizar la reproducibilidad de los experimentos:
\begin{itemize}
  \item Se fija una semilla aleatoria común.
  \item Se documentan las versiones de las librerías (\texttt{requirements.txt}).
  \item La partición de datos por paciente se guarda en disco para reutilizarla.
  \item Cada experimento se ejecuta desde un \textit{notebook} autocontenido.
\end{itemize}

\section{Visualización}
Se desarrollan visualizaciones específicas para el análisis topológico:
diagramas de persistencia y códigos de barras superpuestos sobre la imagen,
mapas de calor de las imágenes de persistencia, y comparación de las regiones
más informativas frente a las activaciones de la CNN (p.\,ej. Grad-CAM).
